\documentclass[12pt]{article}
\usepackage[a4paper,margin=18mm]{geometry}
\usepackage[T1]{fontenc}
\usepackage{amsmath,amssymb}
\usepackage{enumitem}
\usepackage{hyperref}
\usepackage{xr-hyper}
\usepackage[nameinlink]{cleveref}

% Read labels from the lecture build data and link them to lecture01.pdf.
\externaldocument{.build/lecture01}[lecture01.pdf]

\setlength{\parindent}{0pt}
\pagestyle{empty}

\begin{document}

\begin{center}
  {\LARGE\bfseries Tutorial 1: Number Systems}\\[0.25em]
  {\large AMA3707 Real Analysis}\\[0.15em]
\end{center}

The definition references below are imported from the lecture notes. Click a
reference to open the corresponding definition in \texttt{lecture01.pdf}.

\begin{enumerate}[label=\textbf{Question \arabic*.},leftmargin=*]
  \item Using \cref{def:natural-numbers}, determine which of
  \[
    -3,\quad 0,\quad 1,\quad 7,\quad \frac{3}{2}
  \]
  are natural numbers.

  \item Refer to \cref{def:integers}. Decide whether each of
  \[
    -5,\quad 0,\quad 2.4,\quad \sqrt{9}
  \]
  is an integer, and briefly justify each answer.

  \item Using \cref{def:rational-numbers}, prove that \(0.125\) is
  rational by writing it in the form \(p/q\), where \(p,q\in\mathbb{Z}\) and
  \(q\neq0\).

  \item Use \cref{def:integers,def:rational-numbers} to explain
  why every integer is also a rational number.
\end{enumerate}

\end{document}

